\section{Setting and Definitions}\label{sec:setting}

We specify the equations, admission rule and frozen affine model before comparing
regional uncertainty under changes of scale and predictive centre.

\emph{Equations and trained networks.}
On $(x,t)\in[-1,1]\times[0,1]$, the viscous Burgers equation is
\[
 u_t+u u_x=\nu u_{xx},\qquad \nu=0.01/\pi,\qquad
 u(x,0)=-\sin(\pi x),\qquad u(\pm1,t)=0.
\]
The positive-viscosity solution remains smooth and develops a thin layer near $x=0$.
We use the inviscid breaking time $1/\pi$ to define pre- and post-formation windows.
The Allen--Cahn equation is
\[
 u_t=\varepsilon^2u_{xx}+5(u-u^3),\qquad \varepsilon^2=10^{-4},\qquad
 u(x,0)=x^2\cos(\pi x),
\]
with periodic boundaries and interfaces of width $O(\varepsilon)$ between phases
$u\approx\pm1$. A contextual linear control is recorded in Appendix~\ref{app:training};
it is outside the nonlinear scale study.

The independent-feature MLP has $20$ spatial and $20$ temporal sub-networks feeding
a fusion head, with $P=196{,}009$ parameters. Replacing its input layer by random
Fourier features gives $P=275{,}369$ at bandwidth $\mathrm{FS}$
\citep{tancik2020fourier}. The parameter vector $\theta\in\mathbb R^P$ contains
all trainable network weights and biases. The sampled Fourier frequencies are
held fixed throughout training and uncertainty analysis and are excluded from
$\theta$. Both networks are trained in float64 by Adam
\citep{kingma2015adam} followed by Levenberg--Marquardt
\citep{levenberg1944,marquardt1963}. Evaluation uses a uniform $400\times200$ grid ($80{,}000$
points); the residual Jacobian has $N=6400$ rows: $5000$ collocation, $1000$
initial-condition and $400$ boundary rows. Training locations are sampled uniformly
once per seed and held fixed; seeds jointly vary these locations, network
initialisation and, for Fourier models, embedding frequencies. The Burgers reference uses the Cole--Hopf
solution with $80$-node Gauss--Hermite quadrature, agreeing with independent
log-space quadrature to $4.5\times10^{-14}$. The Allen--Cahn reference uses
fourth-order exponential time differencing with $512$ modes and $1000$ steps
\citep{kassam2005etdrk4}; refinement to $2048$ modes changes the canonical contrast
by less than $0.1\%$. Appendix~\ref{app:training} records the training details.
Unless a seed count is stated, numerical illustrations use Burgers--MLP seed~$2$
or Allen--Cahn seed~$0$.

\begin{samepage}
\begin{definition}[Region partition]\label{def:regions}
Let $p_{70},p_{95}$ be the $70$th and $95$th percentiles of $|u_x^\ast|$ over the
whole evaluation grid. Define
\[
 \mathcal S=\{|u_x^\ast|\ge p_{95}\},\qquad
 \mathcal B=\{|u_x^\ast|<p_{70}\},\qquad
 \mathcal T=(\mathcal S\cup\mathcal B)^c.
\]
These are the sharp, bulk and transitional regions.
\end{definition}
\end{samepage}
The reference partition enters admission (Definition~\ref{def:gate}), evaluation and
reference-informed diagnostics; it is absent from training and covariance construction.
It follows moving and narrowing
features: the Allen--Cahn interface moves by $58$ grid cells, and sharp Burgers cells
per slice fall from $58$ at $t=0.13$ to $8$ at $t=0.5$. No sampled reference-gradient
magnitude equals either threshold. The top-$5\%$ definition alone does not establish resolution
(Appendix~\ref{app:grid}). For Burgers, the sharp threshold is $3.34$, against
maximum $116.8$, and $69\%$ of sharp points precede formation; temporal results
therefore also use the two windows separately.

\begin{samepage}
\begin{definition}[Run admission]\label{def:gate}
A trained prediction is admitted if its relative error outside the sharp region satisfies
\begin{equation*}
 \frac{\| (u_\theta-u^\ast)|_{\mathcal B\cup\mathcal T}\|_2}
      {\| u^\ast|_{\mathcal B\cup\mathcal T}\|_2}<\tau,
 \qquad \tau=0.03\text{ for Allen--Cahn},\quad0.02\text{ otherwise}.
\end{equation*}
\end{definition}
\end{samepage}
These trained-network decisions remain fixed in every scale and centre comparison.
Excluding $\mathcal S$ avoids selecting directly on sharp error, but selecting a small
bulk error can enlarge its regional ratio. Burgers--MLP has no seed filtering
($8/8$ admitted); the other nonlinear families are selected samples. A common
$2\%$ threshold retains $13/17$ runs; Appendix~\ref{app:training} reports the
saved-output comparison on that subset.

\begin{table}[t]
\centering\small

\begin{tabular}{@{}lrrrr@{}}
\toprule
configuration & admitted/attempted & error range & median & $\tau$\\
\midrule
Fisher--KPP & $0/5$ & $81.7$--$170.5\%$ & $100.0\%$ & $2\%$\\
Nagumo, both initial conditions & $0/10$ & $75.9$--$90.3\%$ & $80.5\%$ & $2\%$\\
Kuramoto--Sivashinsky, $T=100$ & $0/5$ & $98.2$--$103.8\%$ & $100.4\%$ & $2\%$\\
Kuramoto--Sivashinsky, $T=10$ & $0/5$ & $9.9$--$17.0\%$ & $12.8\%$ & $2\%$\\
\midrule
Burgers, MLP & $8/8$ & $0.05$--$0.40\%$ & $0.08\%$ & $2\%$\\
Burgers, Fourier & $3/8$ & $0.37$--$1.48\%$ & $0.78\%$ & $2\%$\\
Allen--Cahn, MLP & $3/5$ & $1.39$--$2.33\%$ & $2.33\%$ & $3\%$\\
Allen--Cahn, Fourier & $3/4$ & $1.64$--$2.60\%$ & $2.34\%$ & $3\%$\\
Convection--diffusion, MLP control & $3/4$ & $0.19$--$0.31\%$ & $0.21\%$ & $2\%$\\
\bottomrule
\end{tabular}
\caption{Admission accounting and relative $L^2$ error on $\mathcal B\cup\mathcal T$.
Rejected rows include every attempted run; admitted rows include admitted runs only.
Nagumo pools five seeds for each of two initial conditions. The control varies paired
pulse width and diffusivity at one seed. Its narrowest MLP run is rejected at
$62.2\%$; a separate Fourier attempt is rejected at $63.3\%$. All Kuramoto--Sivashinsky
attempts also fail a relaxed $5\%$ gate. Allen--Cahn MLP has median $2.330\%$
and maximum $2.331\%$; both round to $2.33\%$.}
\label{tab:rejected}
\end{table}

Table~\ref{tab:rejected} includes every attempted setting. A small training residual
does not ensure an accurate field: one Kuramoto--Sivashinsky run has loss
$4.4\times10^{-5}$ at $103.8\%$ relative error on $\mathcal B\cup\mathcal T$. The rejected settings therefore
measure failure of this training-and-collocation protocol, consistent with known
stiff-problem difficulties \citep{krishnapriyan2021failure,wang2022ntk}.

\emph{Working likelihood, prior and centres.}
Write $r(\theta)\in\mathbb R^N$ for the stacked training residuals
$r_{\rm blk}=w_{\rm blk}e_{\rm blk}$. The fixed factors are $w=10$ on initial conditions,
$w=10$ on Burgers Dirichlet boundaries or Allen--Cahn periodic value matching,
and $w=1$ on collocation and periodic derivative matching. With common residual
variance $v$, the working likelihood is $r(\vartheta)\sim\mathcal N(0,vI)$; hence
the baseline negative log-likelihood is
\begin{equation*}
 -\log p=\tfrac12\|r(\vartheta)\|^2+\mathrm{const}
 =\tfrac12\sum_{\rm blk}w_{\rm blk}^2\|e_{\rm blk}(\vartheta)\|^2+\mathrm{const}
 \qquad(v=1).
\end{equation*}
The weights specify the likelihood as well as the training loss: the unweighted
block noise standard deviations are $\sqrt v/w_{\rm blk}$.
All analyses retain this raw weighted metric, with $J=\partial r/\partial\theta$
and $F=J^\top J$. For $\alpha,v>0$, the isotropic prior is $\vartheta\sim\mathcal N(0,\alpha^{-1}I)$.

\begin{samepage}
The baseline uses $v=1$ and $\alpha_0=1/\langle\theta^2\rangle$, where
$\langle\theta^2\rangle=P^{-1}\|\theta\|^2$ averages over all trainable weights
and biases; this plug-in rule is not evidence maximisation.
\end{samepage}

Freeze the affine residual $r(\theta)+J(\vartheta-\theta)$ and affine output
$u_\theta(x)+g(x)^\top(\vartheta-\theta)$, where $g=\nabla_\theta u_\theta$.
In predictive expressions, $x$ denotes the full space--time input $z$ of
Section~\ref{sec:intro}; in the PDEs it denotes the spatial coordinate.
Conditioning on zero residual gives
\[
 \Sigma(\alpha,v)=v(F+\gamma I)^{-1},\qquad
 \sigma^2(x)=g(x)^\top\Sigma g(x),\qquad \gamma=\alpha v,
\]
and the exact affine posterior mean
\begin{equation}\label{eq:affine-mean}
 \delta_\gamma=-(F+\gamma I)^{-1}[J^\top r(\theta)+\gamma\theta],
 \qquad m_\gamma(x)=u_\theta(x)+g(x)^\top\delta_\gamma.
\end{equation}
The prior is absent during training, so $u_\theta$ need not equal $m_\gamma$ even
at a stationary training solution. A $\theta$-centred (checkpoint-centred) prior
$\vartheta\sim\mathcal N(\theta,\alpha^{-1}I)$ instead gives
\[
 \delta_{L,\gamma}=-(F+\gamma I)^{-1}J^\top r(\theta),\qquad
 m^\theta_\gamma=u_\theta+g^\top\delta_{L,\gamma}.
\]
Its covariance is unchanged. This prior is conditional on the fitted checkpoint.
We compare its affine mean at the stated scales and available trained-centre roots,
and evaluate $\widetilde u^\theta_\gamma=u_{\theta+\delta_{L,\gamma}}$ as a nonlinear
forward diagnostic. The covariance remains at $\theta$; this diagnostic is not a
nonlinear posterior predictive mean. Prediction bounds are in
Appendix~\ref{app:r3:centres}.
At fixed $\gamma$, centres and uncertainty contrasts are unchanged by matched
$(\alpha,v)$ choices, while widths scale as $\sqrt v$. No output-noise term is added
when comparing latent uncertainty with the deterministic reference.

Baseline variances use the Woodbury identity
\begin{equation*}
 \sigma^2(x)=\alpha^{-1}\left[\|g\|^2-(Jg)^\top
 (JJ^\top+\alpha I)^{-1}Jg\right].
\end{equation*}
The recorded baseline solves required no added diagonal regularisation. Weak-precision
calculations instead factor $J$ directly by QR/SVD to avoid cancellation
(Appendix~\ref{app:n1}).

\emph{Regional diagnostics.}
For each predictive centre, write $e$ for prediction minus $u^\ast$ and
$\langle\cdot\rangle_r$ for a regional mean.
\begin{samepage}
\begin{definition}[Contrast, target and gap]\label{def:metrics}
\[
 \Rrms=\sqrt{\frac{\langle\sigma^2\rangle_{\mathcal S}}
                         {\langle\sigma^2\rangle_{\mathcal B}}},\qquad
 \Rtar=\sqrt{\frac{\langle e^2\rangle_{\mathcal S}}
                        {\langle e^2\rangle_{\mathcal B}}},\qquad
 \ratiogap=\Rtar/\Rrms.
\]
The secondary mean-width contrast is
$\achratio=\langle\sigma\rangle_{\mathcal S}/\langle\sigma\rangle_{\mathcal B}$.
\end{definition}
\end{samepage}
Ratios assume positive denominators.
At fixed scales, all centres share the contrast; the target, gap and coverage differ.
The identity $\ratiogap=1$ matches regional RMS contrasts, not pointwise calibration
or absolute width. We also report
\begin{equation*}
 \meanZ_r=\langle|e|/\sigma\rangle_r,
 \qquad \mathrm{cov}@k_r=\langle\mathbf1\{|e|/\sigma\le z_k\}\rangle_r,
\end{equation*}
with $z_{68}=1$, $z_{95}=1.96$. Coverage is the fraction of grid points whose intervals contain the deterministic
reference, not a repeated-data coverage estimate;
grid points are not independent replicates. Under a correctly specified centred
Gaussian error model the expectations are $\sqrt{2/\pi}\approx0.8$ and approximately $0.68,0.95$.
Throughout, sampled curves describe evaluated settings; refined roots locate the
observed trained-centre crossings within retained numerical brackets. Direct root
fields are compared with the coverage envelopes.
A \emph{best tested evidence setting} maximises a finite search, and a
\emph{consistency check} compares numerical representations. All spectral conclusions
use the returned affine model; the checks do not certify weak-mode forward accuracy.
